% ==================== 3 技术选择：优先红外还是拉曼 ====================
\section{技术选择}

\subsection{优先红外光谱}

在技术选择上，红外具有明显优势，原因如下：1. 判据直接。COO$^-$ 的 $\nu_{\mathrm{as}}$ 与 $\nu_{\mathrm{s}}$ 在红外里都是强峰，拉曼截面却小得多。2. 判据可量化。配位方式由 $\Delta\nu=\nu_{\mathrm{as}}-\nu_{\mathrm{s}}$ 判定：明显小于 110 cm$^{-1}$ 为双齿螯合，接近或大于 200 cm$^{-1}$ 为单齿，桥联介于 140--200 cm$^{-1}$\cite{deacon1980}。3. 干扰可识别。1710 cm$^{-1}$ 的 C=O 若仍然很强，表明样品里还有游离油酸\cite{piao2026ir}。与油酸、油酸钠及模型化合物做差谱与分峰，三种配位可区分（表~\ref{tab:compare}）。

\begin{table}[H]
  \centering
  \caption{红外与拉曼的比较}
  \label{tab:compare}
  \small
  \begin{tabularx}{\textwidth}{@{}lXX@{}}
    \toprule
    项目 & 红外（ATR-FTIR） & 拉曼（785 nm） \\
    \midrule
    物理过程 & 吸收，$\Delta\mu\neq0$ & 非弹性散射，$\Delta\alpha\neq0$ \\
    本体系关键信息 & COO$^-$ 的 $\nu_{\mathrm{as}}$/$\nu_{\mathrm{s}}$、$\Delta\nu$、C=O、O—H & 144/447 cm$^{-1}$、C—H/C=C \\
    水与干扰 & 水干扰强；KBr 吸潮 & 水峰弱；荧光强、信号弱 \\
    空间分辨率 & 1.5--10 $\mu$m & 250--500 nm \\
    \bottomrule
  \end{tabularx}
\end{table}

\subsection{拉曼光谱能够补充的信息}

拉曼光谱能够补充以下信息：1. 晶相组成。锐钛矿 144 cm$^{-1}$ 与金红石 447 cm$^{-1}$ 峰形尖锐，峰强比或 PLS 模型都能定量，5\%--95\% 内偏差可优于 2\%\cite{sacco2023}，红外低波数吸收难以做到。2. 空间分布。油酸的 C=C 与 C—H 在拉曼里也有响应，样品可在水相直接测量并做 mapping\cite{kurouski2020}。不足是 COO$^-$ 信号弱，$\Delta\nu$ 的判据不如红外可靠。
